% ==================================================================
% id:       cond-018
% topic:    הסתברות מותנית ואי-תלות
% subtopic: נוסחת ההסתברות השלמה ונוסחת בייס
% solution: full
% ==================================================================
בארץ שני מפעלים למגבות: מפעל $A$ מייצר 60\% מהמגבות, ומפעל $B$ מייצר 40\%. מחצית מהמגבות של $A$ הן מסוג א׳ ומחציתן מסוג ב׳. במפעל $B$ רק 20\% מהמגבות הן מסוג ב׳. לחו״ל מייצאים רק מגבות מסוג א׳: מפעל $A$ מייצא 30\% מהמגבות שלו מסוג א׳, ומפעל $B$ מייצא 50\% מהמגבות שלו מסוג א׳.
\begin{enumerate}
\item כמה אחוזים מהמגבות מיוצאים לחו״ל?
\hint{מגבת מיוצאת רק אם היא מסוג א׳. חשבו קודם, לכל מפעל, את ההסתברות שמגבת ממנו מיוצאת.}
\begin{solbox}
נסמן $E$ = ״המגבת מיוצאת״, $T$ = ״המגבת מסוג א׳״. לפי כלל הכפל
\[
\PP(E\mid A)=\PP(T\mid A)\,\PP(E\mid A\cap T)=0.5\cdot0.3=0.15,\qquad
\PP(E\mid B)=0.8\cdot0.5=0.4.
\]
לפי נוסחת ההסתברות השלמה
\[
\PP(E)=0.15\cdot0.6+0.4\cdot0.4=0.09+0.16=\ans{0.25}.
\]
\end{solbox}
\item מגבת שנבחרה באקראי מיוצאת לחו״ל. מה ההסתברות שהיא ממפעל $B$?
\begin{solbox}
לפי בייס, $\PP(B\mid E)=\frac{0.16}{0.25}=\ans{0.64}$.
\end{solbox}
\end{enumerate}
